\color{unidarkgreen}{\subsection{Контроль цепей напряжения (КЦН)}}
\color{black}

\begin{enumerate}[label=\arabic{section}.\arabic{subsection}.\arabic*,  align=left, labelsep=4pt, leftmargin=0pt, itemindent=57pt]

\item
Контроль цепей напряжения предназначен для выявления нарушений в цепях напряжения низкой и средней сторон трансформатора.

\item
Алгоритм работы КЦН приведен на рисунке \ref{pict:KCN}. Уставки КЦН приведены в таблице \ref{app_set:kcn}.

\begin{figure}[H]
\centering
\includegraphics[width=0.9\textwidth,height=0.9\textheight,keepaspectratio]{\fbpath/07.01 РЗТ/081.1741 КЦН/img/KCN.pdf}
\captionof{figure}{Функциональная схема логики КЦН}\label{pict:KCN}
\end{figure}

\item
Цепи напряжения подключаются к ТН, установленным на шинах низшего и среднего напряжения.

\item
Контроль цепей напряжения низкой и средней стороны трансформатора основан на выявлении нарушения симметрии в цепях напряжения и просадки линейного напряжения соответствующей стороны. 

\item
Ввод в работу КЦН осуществляется независимо для каждой стороны с помощью уставок \textbf{«КЦН НН1(НН) 1СШ»}, \textbf{«КЦН НН1(НН) 2СШ»}, \textbf{«КЦН НН2(СН) 1СШ»} и \textbf{«КЦН НН2(СН) 2СШ»}. 

\item
Порог срабатывания минимального линейного напряжения задается уставками \textbf{«Uл КЦН НН1(НН) 1СШ»}, \textbf{«Uл КЦН НН1(НН) 2СШ»}, \textbf{«Uл КЦН НН2(СН) 1СШ»} и \textbf{«Uл КЦН НН2(СН) 2СШ»}. 

\item
Порог срабатывания максимального напряжения обратной последовательности задается уставками \textbf{«U2 КЦН НН1(НН) 1СШ»}, \textbf{«U2 КЦН НН1(НН) 2СШ»}, \textbf{«U2 КЦН НН2(СН) 1СШ»} и \textbf{«U2 КЦН НН2 СН) 2СШ»}. 

\item
Выдержка времени контроля неисправности цепей напряжения регулируется уставками \textbf{«Тср КЦН ТН НН1(НН) 1СШ»}, \textbf{«Тср КЦН ТН НН1(НН) 2СШ»}, \textbf{«Тср КЦН ТН НН2(СН) 1СШ»} и \textbf{«Тср КЦН ТН НН2(СН) 2СШ»}.

\item
Оперативный ввод/вывод КЦН производится с помощью БУ (блоков управления функциями) \textit{«Вывод КЦН НН1(НН) 1СШ»}, \textit{«Вывод КЦН НН1(НН) 2СШ»}, \textit{«Вывод КЦН НН2(СН) 1СШ»} и \textit{«Вывод КЦН НН2(СН) 2СШ»}.


\end{enumerate} 